\scanpage{70}
\mspage{E.}
% archive foliation: 69
Le dernier \unsure{est} est appelé le \uline{groupe fondamental birationnel}
de $K/k$.
[N.B. il y aurait à l'introduire plutôt comme groupe alg. \unsure{lié|lin}
sur $K$],
\inserted{noté $\pi_1(K/k;\Omega)$.}
\inserted{Pour tt modèle propre \uline{projectif} $X$,
$\pi_1(X;\Omega)$ est un quotient de $\pi_1(K/k;\Omega)$.}
On voit donc qu'il est, de plus, un quotient de $\pi_1(U)$.
Donc, comme
\inserted{$U$ peut être pris complémentaire d'une partie fermée de codim $\geq2$
dans un schéma intègre propre;}
\uline{Th.} Si $K/k$ est une extension de type fini avec $k$ alg. clos,
alors
\deleted{\unsure{Alors}}
$\pi_1(K/k;\Omega)$ est top$ {}^t$ de type fini.

\margin{\uline{Remarque}
Si $K$ a un modèle propre régulier, alors pour $X'$ propre normal
dominant $X$, on a
$\pi_1(X;\Omega)\simeq\pi_1(X';\Omega)$,
donc le système projectif des $\pi_1(X_i;\Omega)$,
$X_i$ propre normal, est essentiellement constant.
Qu'en est-il en général ??? Si $X$ est un modèle propre normal,
\deleted{\unsure{tout}}
$U$ l'ouv de ses pts réguliers, est-il vrai que
$\pi_1(U;\Omega)\to\pi_1(K/k;\Omega)$ est bijectif,
i.e. que tte
\deleted{\unsure{modèle $X$}}
extension finie $K'$ de $K$,
non ramifiée sur $U$, est non ramifiée
sur un modèle $X'$ \unsure{déduit de}
$X$ \unsure{par éclatemts} ?
C'est une question au fond de \unsure{schémas génériques} \ldots}
